% 1. La discussion sur les possibilités d'utilisation de nouveaux capteurs pour mesurer la
% durée d'ensoleillement et l'anémomètre

\subsection{Ensoleillement}
Plusieurs choix ont été analysé. Cependant, du moment qu'un capteur précis est nécessaire, les prix monte jusqu'à atteindre la lune.
Le capteur actuel ne peut simplement pas être utilisé. Il n'a pas les données nécessaires pour donnée une unité d'ensoleillement
qui est fidèle.
Le tableau suivant présente une analyse simple de capteur commerciaux trouvé.

\begin{table}[h]
\centering
\small
\resizebox{\textwidth}{!}{%
\begin{tabular}{|l|c|c|c|}
\hline
\textbf{Caractéristique} & \textbf{LS06-S} & \textbf{Apogee SP-215-SS} & \textbf{Kipp \& Zonen SMP10} \\
\hline
Prix & $<5\$ $ & $\approx 339\$ $ & $\approx 1500\$ $ \\
\hline
Précision & Non spécifiée & $<3\%$ à $1000\ \mathrm{W/m^2}$ & ISO 9060 Classe A à réponse spectrale plate \\
\hline
Résolution & Non spécifiée & Dépend du CAN & $0.1\ \mathrm{W/m^2}$ \\
\hline
Fiabilité & Élevée & Calibration usine, dérive $<2\%$/an & Extrême \\
\hline
Fidélité & Non spécifiée & Répétabilité $<1\%$ & Extrême \\
\hline
Tension d'alimentation & 5V & 5.5 à 24V DC & 5 à 30V DC \\
\hline
Courant & 18--30\,\mu A & 300\,\mu A & 8mA \\
\hline
Protocole & Analogique & Analogique 0--5V & UART \\
\hline
Plage spectrale & Visible seulement & 360--1120\,nm & 285--3000\,nm \\
\hline
Temps de réponse & Non spécifié & $<1$ ms & $<1$ s \\
\hline
Température de fonctionnement & Non spécifiée & $-40$ à $70^\circ$C & $-40$ à $80^\circ$C \\
\hline
Plage d'humidité & Non spécifiée & 0--100\% RH & 0--100\% RH \\
\hline
Champ de vision & Non spécifiée & $180^\circ$ & $180^\circ$ \\
\hline
Type de sortie & Analogique & 0--5V & Numérique / analogique \\
\hline
Certification & Aucune & ISO 9060:2018 Classe C & ISO 9060 Classe A \\
\hline
Dérive annuelle & Non spécifiée & $<2\%$/an & Très faible \\
\hline
Résistance à l'eau & Non spécifiée & Submersible jusqu'à 30m & IP67 \\
\hline
Garantie & Non spécifiée & 4 ans & 5 ans typique \\
\hline
\end{tabular}%
}
\caption{Comparaison des options de pyranomètres}
\label{tab:pyranometer_comparison}
\end{table}

Le capteur à 1500\$ est évidament le meilleur possible. Mais il est beaucoup trop chère et son niveau scientifique de précision n'est
pas requis. Le SP-215-SS est la solution idéal. Il peut être complètement éteint lorsqu'il n'est pas utilisé, à une consommation extrèment
faible, est conçus pour des environnement complex, requiert aucun protocol outre une lecture analogue, et est même calibré selon des normes ISO.

Les données lu par ce capteur sont fiable et utilisable. De plus, la fiche de donnée fournis par la compagnie est complète, précise et contient
toutes les données nécessaires pour les acheteurs de la station météo. C'est un capteur de qualité industriel. Il sera donc choisi malgrés sont
prix. Les utilisateurs pourront l'inclure ou non dans leurs configuration lors de l'achat de la sation météo.

\subsection{L'anémomètre}
La même étude que le capteur d'ensoleillement à été faite, et la même conclusion à été trouvé.
Le capteur actuel ne peut simplement pas être utilisé. Il ne résistera pas aux vent et à l'hivers, est fait de plastique qui ne résiste pas
au soleil, et est aussi fragile qu'une tige de spaghetti.
Le tableau suivant présente les alternatives les plus intérresant qui ont été trouvé.

Il est important de se rappellé que l'unité SparkFun contient 3 senseurs.

\begin{table}[h]
\centering
\small
\begin{tabular}{|l|c|c|c|}
\hline
\textbf{Caractéristique} & \textbf{SparkFun} & \textbf{Inspeed WS2R} & \textbf{Gill WindSonic 60} \\
\hline
Prix & 80\$ & 96\$ & 600\$ \\
\hline
Précision & Inconnue & $\pm 4\%$ & 2\% à 12 m/s \\
\hline
Résolution & $\approx 2.4$ km/h à non spécifié & $\approx 1.6$ km/h à $>200$ km/h & 0.01 m/s, 0.1° à 75 m/s \\
\hline
Fiabilité & Très pauvre & Correct (calibré) & Extrême (IP66) \\
\hline
Fidélité & Correct & Meilleur & Extrême \\
\hline
Tension (V) & Source externe & Source externe & 12--30 VDC \\
\hline
Courant & Dépend de la direction & 3 mA (capteur Hall) & 2 mA à 44 mA \\
\hline
Protocole & Lecture analogique & Lecture analogique & RS-232 \\
\hline
\end{tabular}
\caption{Comparaison de capteurs de vent}
\label{tab:wind_sensors}
\end{table}

Le meilleur capteur est le Gill WindSonic 60. Il à une extrème précision, dit la direction du vent en degrées à une extrème précision.
Il contient tout le nécessaire dans ça feuille de données. De plus, son gros avantage est qu'il n'a aucune pièce qui bouge, fessant en
sorte qu'il est garanti 5 ans par la compagnie. Ils ont même certains modèles plus couteux qui ont des fonctionalitées pour enlevé la glace
et la neige. Le ESP32 supporte le protocole RS-232 nativement, donc communiqué avec est simple.

Il vas avoir besoin d'un circuit pour le 12V nécessaire, mais il à juste besoin d'être actif lors des mesures. La calibration et la précision
pour son prix, malgrés qu'il est très chère, justifie son addition à une station météo robuste et durable.


